\chapter[A company reconstructed and stressed]{A complete small company: from reports to divergent stress outcomes}
\label{sim:company-case}

The earlier examples isolate individual mechanisms. This chapter joins them in one synthetic company so that every projected balance can be traced to an event. The company is deliberately simple: one currency, one legal entity, no tax, no dividends, no acquisitions and no other comprehensive income. These exclusions make the arithmetic inspectable. They do not describe a complete real-company model. The point that carries over is the need to preserve compatible economic histories and their joint statement consequences.

\section{The opening statement}
At the beginning of the historical year, the company has cash 40, receivables 30, inventory 50 and property carrying amount 80. Its liabilities are debt 100 and trade payables 40. Equity is 60. All amounts use one currency unit.
\begin{center}
\begin{tabular}{lr|lr}
\toprule Assets&Amount&Liabilities and equity&Amount\\\midrule
Cash&40&Debt&100\\
Receivables&30&Trade payables&40\\
Inventory&50&Equity&60\\
Property&80&&\\\midrule
Total&200&Total&200\\\bottomrule
\end{tabular}
\end{center}
The equality is an accounting starting point. It does not identify customers, inventory units, asset vintages or debt terms. We initially admit the balances as exact to focus on hidden composition. A later section reintroduces rounding and source uncertainty.

The company operates during the year and reports sales 120, cost of sales 70, cash operating expenses 25, interest expense 5 and depreciation 10. With no tax in this teaching model, profit is 10. The cash-flow information reports customer receipts 130, supplier payments 50, operating-expense payments 25, interest payments 5 and capital expenditure 20. No debt is issued or repaid.

\section{One compatible event history}
One history divides sales into cash sales 80 and credit sales 40. Collections on receivables total 50: 20 from opening receivables and 30 from current-year credit sales. Customer receipts therefore total $80+50=130$. Purchases on credit total 60. Supplier payments are 50.

The receivables and inventory bridges are
\begin{align}
 A_1&=30+40-50=20,\\
 I_1&=50+60-70=40.
\end{align}
Trade payables become $40+60-50=50$. Property becomes $80+20-10=90$. Cash becomes
\begin{equation}
 C_1=40+130-50-25-5-20=70.
 \label{sim:company-cash}
\end{equation}
Profit of 10 increases equity from 60 to 70. Thus closing assets are $70+20+40+90=220$, equal to debt 100, payables 50 and equity 70. Every amount follows from a stated movement.

This is one witness. The reports may not disclose the split of collections between opening and new receivables, and they do not disclose which customers remain unpaid. Other event histories can generate the same statements. The model should not turn the chosen witness into an assertion about what actually happened.

\section{The journal bundle}
Under the illustrative recognition rules, cash sales debit cash and credit revenue by 80. Credit sales debit receivables and credit revenue by 40. Collections debit cash and credit receivables by 50. Purchases debit inventory and credit payables by 60. Consumption of inventory debits cost of sales and credits inventory by 70.

Supplier payments debit payables and credit cash by 50. Operating expenses debit expense and credit cash by 25. Interest payment debits interest expense and credits cash by 5. Capital expenditure debits property and credits cash by 20. Depreciation debits expense and credits accumulated depreciation or the summarized carrying-value account by 10.

Each posting bundle balances. More importantly, the bundles reconcile the economic movements. The inventory purchase is not expensed a second time when cash is paid. Credit sales are not recognized again on collection. Capital expenditure is not simultaneously charged as the full current expense and depreciated. These familiar errors are easy to introduce when separate forecasting formulas are joined without event identities.

The summarized property account hides gross cost and accumulated depreciation. That is acceptable for this bounded calculation, but a later disposal or impairment query may require those components. The representation should retain the option to refine rather than declare the summary universally sufficient.

\section{What the reports leave unknown}
At the forecast origin, receivables of 20 can be concentrated in one customer or spread among ten customers owing 2 each. The debt of 100 can be fixed-rate at five per cent or floating-rate with a historical rate of five per cent. Both alternatives reproduce the historical interest payment and closing debt.

Let $c$ indicate concentration: $c=1$ for one customer and $c=0$ for ten equal customers. Let $f$ indicate floating debt: $f=1$ for floating and $f=0$ for fixed. For this exercise we admit all four combinations because the specified evidence provides no linking restriction. This is an explicitly declared feasible family, not a claim that real financing structure and customer concentration are statistically independent. No probabilities are assigned.

Other hidden quantities remain fixed in this first stress calculation. Inventory and property do not change during the short stress horizon, and there are no new sales or purchases. This isolates collection and financing exposure. An operating forecast would add production, demand and replenishment policies rather than assume these balances remain static indefinitely.

\section{Define the shock precisely}
At the start of the stress horizon, payables of 50 must be paid. Next, one largest customer (any one of the tied customers in the diversified case) defaults and pays nothing on its outstanding receivable; every other customer pays in full. At the end of the horizon, one full year's interest is paid after market rates rise by two percentage points. There is no hedge, cap, floor or principal maturity during the horizon.

The customer loss is 20 in the concentrated history and 2 in the diversified history. Write
\begin{equation}
 L(c)=2+18c.
\end{equation}
Collections are $20-L(c)$. Interest is $5+2f$. After paying suppliers, collecting surviving receivables and paying interest, cash is
\begin{align}
 C_2(c,f)&=70-50+[20-L(c)]-(5+2f)\\
 &=33-18c-2f.
 \label{sim:company-stress-cash}
\end{align}
The formula preserves the joint dependence on concentration and rate type. It is an exact result for the declared event sequence.

\section{The four outcomes}
The resulting cash and equity are
\begin{center}
\begin{tabular}{llrr}
\toprule Receivables&Debt&Closing cash&Closing equity\\\midrule
Diversified&Fixed&33&63\\
Diversified&Floating&31&61\\
Concentrated&Fixed&15&45\\
Concentrated&Floating&13&43\\\bottomrule
\end{tabular}
\end{center}
Equity falls by customer loss and interest, so $E_2=70-L(c)-(5+2f)=63-18c-2f$. Receivables and payables both become zero. Inventory remains 40 and property 90. Each history satisfies
\begin{equation}
 C_2+40+90=100+E_2.
\end{equation}
The accounting relation is shared, but the economic outcomes differ by 20 across the family. A single historical witness could conceal that range.

The worst cash outcome is concentrated floating debt. The best is diversified fixed debt. These are feasible joint witnesses, not independent marginal extrema assembled from different worlds. If later evidence rules out one combination, the family and its extrema should update accordingly.

\section{A liquidity threshold changes the question}
Suppose a policy requires closing cash at least 14. Three histories meet it; the concentrated floating history ends at 13. Without probabilities, the correct statement is that the policy is not robustly feasible across the admitted family. It is wrong to say the probability of failure is one quarter merely because there are four rows.

The row count depends on representation. The diversified case could be split into ten different customer-name assignments without changing its economic exposure. Equal weighting would then change the apparent failure probability. A measure must be attached to economic states or events with justified weights, not to a convenient enumeration.

If the threshold is a contractual covenant, its consequences require the actual covenant definition and remedy process. If it is an internal treasury target, a shortfall is a policy deviation rather than legal default. The same numerical comparison can support different meanings depending on the source of the threshold.

\section{A funding response}
Suppose a committed facility can supply up to 2 at the final payment date, with no fee or interest within this short example, and management draws only enough to reach cash 14. The concentrated floating case draws 1. Its closing cash becomes 14 and debt 101; equity remains 43. The balance sheet still reconciles: $14+40+90=101+43$.

The draw improves current liquidity but creates a future obligation. It does not reverse the credit loss or interest expense. A model that raises equity by the draw amount would misclassify financing. A model that treats the draw as free permanent cash would overstate future capacity.

If the facility is unavailable after the customer default, the response is infeasible. If it settles after interest is due, it cannot fund that payment without a bridge. If its covenant depends on equity falling below 44, the worst history may lose eligibility precisely when it needs funding. These are distinct contract variants that should be modeled as conditional cases.

\section{A dependence hidden by aggregate data}
Suppose a private lender report establishes that the floating loan was granted only when the receivables were diversified. Then the concentrated floating combination is excluded. The remaining minimum cash is 15, so the threshold 14 is satisfied by every retained history.

This new conclusion comes from a joint restriction, not from narrower marginal ranges. Concentration still can be either zero or one, and debt still can be fixed or floating. Marginal sets are unchanged. A representation keeping only those marginals cannot express why the worst combination disappeared.

The example demonstrates the value of constraints over named variables. The restriction can be written $f=1\Rightarrow c=0$. A new source should be attached to it, with availability and scope. If the source is later rejected, the excluded combination returns in the revised family.

\section{Learning before the payment}
Suppose management knows the loan type but the public analyst does not. Management can plan a precautionary draw in the floating history without using future information. The analyst's forecast must integrate or bound those conditional management actions. Forcing management to use the analyst's ignorance would understate its information.

Now suppose a customer deterioration signal arrives before supplier payment. A policy may delay discretionary capital spending or seek financing in response. If the signal arrives only after payment, those actions may no longer be available. The date and actor access matter as much as the signal's accuracy.

The current stress example has no discretionary spending. Adding one should identify its benefit, cancellation cost and commitment status. An apparently simple cash-saving policy can be infeasible if the expenditure is already contractually committed. The event sequence prevents that distinction from being lost in an annual forecast.

\section{Historical journal ambiguity beyond contracts}
The historical reports also permit different sales and collection timing within the year. One company may collect most cash early and pay suppliers later; another may pay first and collect near year-end. Annual cash flow and closing cash can agree while minimum cash differs.

For example, take opening cash 40 and aggregate receipts 130 with payments 100 in total. If payments of 100 occur before receipts, a company without credit cannot execute the sequence. If receipts arrive first, it can. The annual closing cash of 70 does not distinguish them. A feasible historical reconstruction must respect intermediate cash and financing events where those are within scope.

This does not mean every annual disclosure allows precise daily reconstruction. It means the model should retain timing uncertainty or state that it cannot answer an intra-year liquidity question. Inventing a uniform monthly schedule gives a point path whose apparent smoothness comes from the interpolation assumption.

\section{Inventory alternatives}
Closing inventory value 40 can represent eight units costing 5 or four units costing 10. If a future order requests six physically interchangeable units, the first history has sufficient stock and the second does not. If products are not interchangeable, even total unit count is inadequate.

A monetary demand shock should therefore not be converted automatically into units using an average historical cost. Selling price, cost, product mix and capacity are different quantities. A future price decline can reduce revenue per unit without changing available physical stock. A supply shock can reduce units while increasing unit cost.

The present company case holds inventory fixed to isolate credit and rates. Extending it requires product or cohort state, procurement lead times and sales rules. The aggregate carrying amount remains a reporting consequence of those states. A richer query justifies refinement of the family rather than invalidating the earlier bounded calculation.

\section{Property and replacement cash}
Property of 90 after depreciation can comprise young assets with long remaining lives or old assets near replacement. Equal historical depreciation does not identify next year's capital needs. A policy that sets capital expenditure equal to depreciation assumes a replacement relation.

Suppose one history has a machine requiring replacement next year for 30 and another can continue for five years with maintenance 2. Their current carrying amounts may agree under different acquisition costs and accumulated depreciation. A shock to financing availability has different operating consequences.

A complete extension should retain capacity and maintenance as well as book values. Replacement can change productivity and costs, so it is not only a cash outflow. The model needs a production or service relation and a management policy. The accounting bridge alone cannot supply them.

\section{Rounding widens the family}
The exact amounts used above make arithmetic transparent. If reported receivables 20 are rounded to the nearest unit, the underlying amount may lie in a reporting interval. Customer loss and collections then become functions over that interval. A cash threshold at 14 can be crossed by values that a point reconstruction misses.

Suppose the concentrated floating case's opening cash is itself rounded from an amount between 69.5 and 70.5 under the declared convention. Its stressed cash lies correspondingly between 12.5 and 13.5 if other terms are exact. A facility of exactly 1 may no longer guarantee reaching 14. The required draw can range from 0.5 to 1.5.

Rounding intervals can also be jointly constrained by a total. Treating each line's rounding error independently may admit totals outside the reported balance-sheet interval. The feasible family should retain the total relation rather than merely widen every number separately.

\section{An explanation the system should produce}
For the worst witness, an explanation should identify the concentrated receivable history, floating debt, customer default, rate increase, supplier payment and lack of an eligible funding response. It should show cash 13 and equity 43 with their bridges. It should identify which facts are observed and which are alternative assumptions.

The explanation should not say that the real company will fail. It should say that a compatible history produces the stated shortfall under the specified shock and policy. If the threshold is internal, it should call it a target shortfall. If a covenant consequence is modeled, it should name that consequence and its premises.

A second explanation can identify the source that would exclude the witness, such as a verified customer concentration note or fixed-rate contract. This turns the simulation into a guide for focused investigation. It also makes disagreement productive: a reviewer can challenge a premise rather than merely distrust the final number.

\section{What this complete case establishes}
The case connects reports, compatible histories, events, journals, cash and equity. It shows why identical statements do not identify shock response, why joint restrictions matter, why financing needs a transaction and why timing matters. All numerical relations are small enough for exact checking.

It does not establish a calibrated corporate model, a universal customer-default process or a reporting-standard treatment for every event. Those would require additional facts and authority. The bounded example is valuable because it exposes the mechanisms without pretending those wider questions are resolved.

An implementation fixture should store the input observations, the four initial history cases, the optional dependence restriction, the event sequence, expected journals and final statements. It should also test the rejected variants: equal branch probabilities without a measure, independent marginal output selection, a late facility draw and a cash-balancing equity plug. These are concrete failures of the target semantics.
